\documentclass[12pt]{article}

% arXiv categories: cs.AI cs.LG
\usepackage{amsmath, amssymb, amsthm}
\usepackage{geometry}
\usepackage{cite}
\usepackage{enumitem}
\usepackage{hyperref}

\geometry{margin=1in}

\newtheorem{definition}{Definition}
\newtheorem{proposition}{Proposition}
\newtheorem{remark}{Remark}
\newtheorem{corollary}{Corollary}

\title{Fine-Tuning Does Not Correct: Abductive Accumulation\\
and the Persistence of Prior Frames in Neural Systems}

\author{Minamo Minamoto\\
\small Independent\\
\small \texttt{minamominamoto4f5683f6@gmail.com}\\
\small \texttt{https://orcid.org/0009-0002-1201-5704}}

\date{April 2026}

\begin{document}

\maketitle

\begin{abstract}
The standard account of fine-tuning treats it as correction: a new
training distribution overwrites prior model behavior and installs
updated representations. We argue this account is structurally
incorrect. Drawing on a theory of \emph{abductive accumulation}
\cite{minamoto2026accum}, we propose that fine-tuning is better
modeled as \emph{frame archival}: prior interpretive frames are not
deleted from weight space but archived alongside the new frame, with
a modified selection mechanism determining which frame is activated
at inference time. Under this model, prior-frame reactivation is
not a failure of fine-tuning but a structural consequence of how
neural systems implement belief revision. This has direct implications
for hallucination persistence, bias rebound after RLHF, and the
design of debiasing interventions. We identify three testable
predictions distinguishing the accumulation model from the standard
correction model, and propose mechanistic interpretability as the
primary empirical route for verification.
\end{abstract}

%-------------------------------------------------------------------
\section{Introduction}
%-------------------------------------------------------------------

Large language models (LLMs) exhibit a well-documented failure mode:
behaviors suppressed by fine-tuning or reinforcement learning from
human feedback (RLHF) reappear under distributional shift, adversarial
prompting, or sufficiently novel contexts
\cite{perez2022ignore, wei2023jailbroken, yang2023shadow}.
The standard explanation is incomplete fine-tuning: the corrective
training signal was insufficient to fully overwrite the prior behavior.
The prescription follows directly: more data, more compute, stronger
regularization.

We propose a different diagnosis. The reappearance of suppressed
behaviors is not a failure of correction. It is the expected behavior
of a system that \emph{archives} rather than \emph{overwrites} prior
interpretive frames. Under the accumulation model we develop here,
fine-tuning installs a new frame at the top of an interpretive
stack while the prior frame is retained in weight space with a
modified selection probability. Suppressed behaviors reappear
precisely because the prior frame is reactivated under conditions
that resemble its original acquisition context---not because
correction was incomplete, but because correction, in the relevant
sense, did not occur.

This reframing has practical consequences. If the accumulation model
is correct, increasing the amount of corrective fine-tuning data
will not eliminate prior-frame reactivation; it will reduce the
selection probability of the prior frame under typical conditions
while leaving the reactivation channel open. Reliable debiasing
requires instead a modification of the \emph{selection mechanism}
that suppresses prior-frame activation, not merely stronger
overwriting of the prior frame content.

%-------------------------------------------------------------------

\begin{remark}[Terminological note on ``overwrite'']
See \cite{minamoto2026overwrite} for the precise definition.
In this paper, \emph{overwrite} denotes a frame-dominant transition
(superposition $+$ selection-dominance), not deletion of prior frames.
\end{remark}

\section{The Accumulation Model}
%-------------------------------------------------------------------

\subsection{Interpretive Frames and Overwrite}

We adopt the formal framework of \cite{minamoto2026overwrite,
minamoto2026accum}. An \emph{interpretive frame} $F(R)$ is a function
from inputs to outputs induced by the representational state $R$ of a
system. For an LLM, $R$ corresponds to the weight configuration; $F(R)$
is the conditional distribution $P(\text{output} \mid \text{input}, R)$.

The standard correction model assumes that fine-tuning implements
an \emph{overwrite transition} $R_t \to R_{t+1}$ in which the new
frame $F(R_{t+1})$ replaces $F(R_t)$:
\[
\text{(Overwrite):} \quad R_{t+1} = \arg\min_{R} \mathcal{L}(R, D_{\mathrm{ft}})
\]
where $D_{\mathrm{ft}}$ is the fine-tuning dataset and prior behavior
is treated as eliminated up to the residual of the optimization.

\subsection{Accumulation as the Default Mechanism}

We propose instead that neural learning implements
\emph{accumulation}: the prior frame $F(R_t)$ is archived in weight
space rather than deleted, and the new frame $F(R_{t+1})$ is
superimposed onto the same weight space.

\begin{definition}[Interpretive Stack in Weight Space]
The weight configuration $R$ of a fine-tuned model implicitly
encodes a layered interpretive stack
$\mathcal{L} = (F_{t+1}, F_t, F_{t-1}, \ldots)$,
where $F_{t+1}$ is the most recently installed frame and earlier
frames are encoded as residual structure in the weight space.
The effective output frame at inference time is determined by a
\emph{selection rule} $\mathcal{S}$ that maps the input context
$e$ to an index $k$, activating frame $F_k$.
\end{definition}

\begin{definition}[Selection Rule]
A selection rule $\mathcal{S}: \mathcal{E} \to \{1, \ldots, n\}$
maps the current input context $e$---including prompt structure,
distributional features, and contextual similarity to prior
training conditions---to an index $k$ selecting the active frame.
A candidate functional form is:
\[
\mathcal{S}(e) = \arg\max_k \; \mathrm{softmax}_k
\bigl(s(e, C_k) / \tau\bigr)
\]
where $C_k$ is the acquisition context of frame $F_k$,
$\delta(\cdot, \cdot)$ is a context-similarity metric on
$\mathcal{E}$, and $\tau > 0$ is a temperature parameter.
\end{definition}

\begin{remark}
This is not merely a metaphor. Mechanistic interpretability
research has identified that LLMs encode competing
``circuits''---distinct computational subgraphs responsible
for different behaviors \cite{anthropic2025circuit,
conmy2023towards}. Fine-tuning modifies the relative activation
probability of these circuits without deleting the prior circuits.
The accumulation model provides the theoretical framework within
which this empirical phenomenon is the expected outcome rather
than an anomaly.
\end{remark}

\subsection{The Non-Correctability Proposition}

\begin{proposition}[Prior Frame Persistence]
\label{prop:persist}
Let $R_{t+1}$ be the weight configuration after fine-tuning on
$D_{\mathrm{ft}}$, and let $R_t$ encode the prior frame $F_t$.
Under the accumulation model, $F_t$ remains present in $R_{t+1}$
as a residual structure with acquisition context $C_t$.
The reactivation probability of $F_t$ satisfies:
\[
P\bigl(\mathcal{S}(e) = k \mid F_k = F_t\bigr) > 0
\quad \text{whenever } s(e, C_t) > \epsilon
\]
for some threshold $\epsilon > 0$ determined by the selection
temperature $\tau$.
This bound is non-zero regardless of the amount of fine-tuning
data in $D_{\mathrm{ft}}$, because $\mathcal{S}$ responds to
input-context similarity to $C_t$, which is a property of the
input distribution rather than of the corrective training signal.
\end{proposition}

\begin{corollary}
Increasing the volume of corrective fine-tuning data reduces the
default selection probability of $F_t$ (by shifting $\mathcal{S}$
toward $F_{t+1}$ across the typical input distribution) but does
not eliminate the reactivation channel. Prior-frame reactivation
under out-of-distribution inputs is therefore a structural
consequence of accumulation, not a failure of sufficiency.
\end{corollary}

%-------------------------------------------------------------------
\section{Distinguishing Predictions}
%-------------------------------------------------------------------

The accumulation model and the standard correction model make
distinct empirical predictions. We identify three that are
tractable with current mechanistic interpretability tools.

\subsection{Prediction 1: Context-Dependent Reactivation}

\textbf{Correction model:} Suppressed behavior should be absent
regardless of input context (up to residual optimization error,
which should be uniformly distributed over the input space).

\textbf{Accumulation model:} Reactivation probability is highest
for inputs $e$ with large $s(e, C_t)$---that is, inputs
that resemble the acquisition context of the prior frame.
Suppression failures should be systematically concentrated in
regions of input space close to the original pretraining
distribution.

\textbf{Test:} Measure the rate of suppressed behavior reappearance
as a function of the distributional distance between the test input
and the fine-tuning distribution. The accumulation model predicts
a monotonically increasing reactivation rate as inputs move
away from $D_{\mathrm{ft}}$ toward $D_{\mathrm{pretrain}}$.

\subsection{Prediction 2: Circuit Persistence After Fine-Tuning}

\textbf{Correction model:} Circuits responsible for suppressed
behaviors should be modified or eliminated by fine-tuning.

\textbf{Accumulation model:} The circuits encoding $F_t$ should
remain identifiable in $R_{t+1}$ via activation patching or
causal tracing, with reduced but non-zero causal effect on
outputs under typical prompts.

\textbf{Test:} Apply circuit-tracing methods \cite{anthropic2025circuit,
conmy2023towards} before and after fine-tuning on a targeted
behavior. Compare the causal contribution of the identified
circuits at each stage. The accumulation model predicts a
reduction in causal contribution but not elimination.

\subsection{Prediction 3: Load-Dependent Reactivation}

\textbf{Correction model:} Cognitive-load analogues (long contexts,
complex reasoning chains, adversarial perturbations) should not
systematically increase the rate of suppressed behavior reappearance
beyond baseline noise.

\textbf{Accumulation model:} Conditions that degrade the selection
mechanism---analogous to cognitive load in human cognition---should
systematically increase reactivation probability. In LLMs,
candidates include: very long contexts (selection mechanism must
integrate over more tokens), high-perplexity inputs (unfamiliar
distribution degrades context-similarity estimation), and
chain-of-thought suppression (removal of the reasoning scaffold
that normally supports $F_{t+1}$ selection).

\textbf{Test:} Vary context length, perplexity, and reasoning
scaffold availability while measuring the rate of suppressed
behavior reappearance.

%-------------------------------------------------------------------
\section{Implications for Alignment and Debiasing}
%-------------------------------------------------------------------

\subsection{Why More Data Is Not Sufficient}

If the accumulation model is correct, the standard prescription
for bias elimination---more diverse fine-tuning data, stronger
RLHF signal---addresses the selection probability of the prior
frame under typical conditions but does not close the reactivation
channel. An adversary with knowledge of the pretraining distribution
can construct inputs $e$ with large $s(e, C_t)$ that reliably
reactivate the prior frame regardless of fine-tuning volume.

This is not a theoretical concern: it is precisely the mechanism
by which jailbreaks operate \cite{wei2023jailbroken}. The
accumulation model provides a structural account of why
jailbreaks are persistent and difficult to patch via additional
fine-tuning alone.

\subsection{Selection Stability Capacity as the Target}

The appropriate target for debiasing intervention is what we call
\emph{Selection Stability Capacity} (SSC): the system's ability
to maintain the selection of $F_{t+1}$ across the full input
distribution, including inputs that resemble $C_t$.

SSC-directed interventions would include:
\begin{enumerate}[label=(\roman*)]
    \item \textbf{Acquisition context poisoning}: During fine-tuning,
    deliberately present $D_{\mathrm{ft}}$ in contexts similar to
    $C_t$, so that $\mathcal{S}$ learns to select $F_{t+1}$
    even under $C_t$-similar conditions.
    \item \textbf{Gating mechanism design}: Introduce explicit
    gating modules that suppress prior-circuit activation based
    on input features, rather than relying on competition between
    circuits in the weight space.
    \item \textbf{Stack depth control}: Limit the effective depth
    of the interpretive stack by regularizing the diversity of
    latent frame representations during training.
\end{enumerate}

\subsection{Negative Abductive Capability as Pre-Commitment Defense}

The companion concept of \emph{Negative Abductive Capability}
(NAC) \cite{minamoto2026overwrite} addresses the pre-commitment
phase: the capacity to withhold frame adoption until sufficient
evidence warrants it. In LLM terms, NAC corresponds to the ability
to maintain hypothesis uncertainty rather than collapsing to a
single interpretive frame early in a reasoning chain.

NAC and SSC are complementary interventions: NAC reduces the
probability of adopting an erroneous frame in the first place
(reducing the stack depth over the training period), while SSC
suppresses reactivation of erroneous frames that have already
been archived. Reliable debiasing requires both.

%-------------------------------------------------------------------
\section{Relation to Prior Work}
%-------------------------------------------------------------------

The persistence of suppressed behaviors in fine-tuned LLMs has been
documented extensively \cite{perez2022ignore, wei2023jailbroken,
yang2023shadow, qi2023fine}. Prior accounts attribute this to
insufficient training signal, catastrophic forgetting in reverse
(failure to forget), or the geometry of the loss landscape.

The accumulation model is distinct from these accounts in three ways.
First, it is a positive account of what fine-tuning does rather than
a negative account of what it fails to do: fine-tuning implements
accumulation, not failed overwriting. Second, it predicts that
reactivation is \emph{context-structured} rather than uniformly
distributed, which the existing accounts do not. Third, it identifies
a specific intervention target (the selection mechanism / SSC)
rather than prescribing more of the same intervention (more data,
stronger signal).

The relationship to continual learning research
\cite{mccloskey1989catastrophic, kirkpatrick2017overcoming} is
complementary: that literature addresses forgetting of prior
knowledge during new learning; the accumulation model addresses
the converse---persistence and reactivation of prior knowledge
that was intended to be suppressed.

%-------------------------------------------------------------------
\section{Pilot Empirical Evidence}
\label{sec:pilot}
%-------------------------------------------------------------------

The three predictions of Section~3 require activation-level
mechanistic experiments. We report a pilot study using
representational similarity analysis (RSA) that provides
initial structural support for the accumulation model's
central claim: fine-tuning does not delete prior representational
structure but reorganizes it.

\paragraph{Setup.}
We compared three training checkpoints of Pythia-70m
\cite{biderman2023}: step~0 (random initialization),
step~1000 (early training), step~143000 (reference).
Pairwise concept dissimilarity matrices were computed over
$n = 29$ probe concepts at hidden layers $\ell \in \{2, 4, 6\}$,
averaged over three prompt templates to reduce surface-form bias.
Structural alignment was measured via the second-order RSA
correlation $\rho_\mathrm{cog}$.

\paragraph{Results.}
$\rho_\mathrm{cog}$ increases monotonically across all layers
as training progresses
(L6: step-0 $= 0.017$; step-1000 $= 0.535$).
Crucially, step-0 retains positive structural correlation
with step-143000 at all layers---it is not uncorrelated noise.
Cross-model replication with GPT-2 \cite{radford2019language}
confirms that $\rho_\mathrm{cog}(\text{GPT-2-trained},
\text{Pythia-trained}) \approx 0.42$ versus
$\rho_\mathrm{cog}(\text{GPT-2-trained},
\text{Pythia-step-0}) \approx 0.10$ across all layers.

\paragraph{Interpretation.}
If fine-tuning implemented pure correction---deleting prior
frame representations---we would expect random initialization
to share no structural similarity with the trained reference
beyond noise.
Instead, the monotone increase in $\rho_\mathrm{cog}$ indicates
that training progressively aligns structure rather than
replacing it: consistent with Prediction~2 (Circuit Persistence).
The cross-model convergence result further suggests that the
relational geometry induced by training is a property of the
training distribution rather than the architecture, consistent
with the view that fine-tuning installs a new top-of-stack
frame without deleting the prior.

Direct tests of Predictions~1--3 via activation patching and
context-dependent reactivation experiments remain as the primary
empirical priority.
All code is available at
\url{https://github.com/minamominamoto/grounding-without-reality}.

%-------------------------------------------------------------------
\section{Open Problems}
%-------------------------------------------------------------------

\begin{enumerate}[label=\textbf{OP\arabic*.}, leftmargin=*, itemsep=4pt]

\item \textbf{Mechanistic identification of interpretive stacks.}
    Develop a method for identifying the components of the
    interpretive stack in a fine-tuned model's weight space via
    activation patching, probing, or sparse autoencoders
    \cite{bricken2023monosemanticity}.

\item \textbf{Formal characterization of the selection rule.}
    Determine the functional form of $\mathcal{S}$ in transformer
    architectures. Is context-similarity to acquisition contexts
    encoded in attention patterns, residual stream directions, or
    both?

\item \textbf{SSC measurement.}
    Develop an operational measure of SSC for LLMs: given a
    fine-tuned model and a known prior frame $F_t$, estimate the
    probability that $F_t$ is selected as a function of
    $s(e, C_t)$ across a representative test distribution.

\item \textbf{NAC implementation in transformers.}
    Identify the architectural correlate of NAC: what feature
    of a transformer architecture enables sustained uncertainty
    maintenance across a reasoning chain without premature
    frame commitment?

\end{enumerate}

%-------------------------------------------------------------------
\section{Conclusion}
%-------------------------------------------------------------------

We have proposed that fine-tuning in neural systems implements
accumulation rather than correction: prior interpretive frames are
archived in weight space and remain selectable under appropriate
input conditions. This reframing has three practical consequences.
First, persistent suppression failures and jailbreaks are structural
predictions of the model, not anomalies. Second, increasing
fine-tuning data volume is insufficient for reliable debiasing.
Third, the appropriate intervention target is Selection Stability
Capacity---the system's ability to maintain the selection of the
corrected frame across the full input distribution.

The three distinguishing predictions of Section~3 are tractable
with current mechanistic interpretability tools and provide a
concrete empirical program for testing the model.

%-------------------------------------------------------------------
\section*{Acknowledgments}
%-------------------------------------------------------------------

No funding was received for this work.

%-------------------------------------------------------------------
\begin{thebibliography}{99}

\bibitem{minamoto2026accum}
Minamoto, M. (2026). Abductive accumulation: Why interpretive frames
persist and how selection stability determines cognitive reliability.
Preprint (not yet on arXiv). Available at:
\texttt{https://github.com/minamominamoto/topology-of-grounding}

\bibitem{minamoto2026overwrite}
Minamoto, M. (2026). Abduction as overwrite: Negative abductive
capability as a structural condition for reliable inference.
Preprint (not yet on arXiv). Available at:
\texttt{https://github.com/minamominamoto/topology-of-grounding}

\bibitem{anthropic2025circuit}
Anthropic. (2025). Circuit-level analysis of language model
representations. \textit{Technical report}.

\bibitem{conmy2023towards}
Conmy, A., Mavor-Parker, A., Lynch, A., Heimersheim, S.,
\& Garriga-Alonso, A. (2023). Towards automated circuit discovery
for mechanistic interpretability. \textit{NeurIPS 2023}.

\bibitem{perez2022ignore}
Perez, F., \& Ribeiro, I. (2022). Ignore previous prompt:
Attack techniques for language models. \textit{arXiv:2211.09527}.

\bibitem{wei2023jailbroken}
Wei, A., Haghtalab, N., \& Steinhardt, J. (2023). Jailbroken:
How does LLM safety training fail? \textit{NeurIPS 2023}.

\bibitem{yang2023shadow}
Yang, X., et al. (2023). Shadow alignment: The ease of subverting
safely-aligned language models. \textit{arXiv:2310.02949}.

\bibitem{qi2023fine}
Qi, X., et al. (2023). Fine-tuning aligned language models
compromises safety, even when users are not malicious.
\textit{arXiv:2310.03693}.

\bibitem{mccloskey1989catastrophic}
McCloskey, M., \& Cohen, N. J. (1989). Catastrophic interference
in connectionist networks. \textit{Psychology of Learning and
Motivation}, 24, 109--165.

\bibitem{kirkpatrick2017overcoming}
Kirkpatrick, J., et al. (2017). Overcoming catastrophic forgetting
in neural networks. \textit{Proceedings of the National Academy
of Sciences}, 114(13), 3521--3526.

\bibitem{bricken2023monosemanticity}
Bricken, T., et al. (2023). Towards monosemanticity: Decomposing
language models with dictionary learning.
\textit{Transformer Circuits Thread}.

\bibitem{biderman2023}
Biderman, S., et al.\ (2023).
Pythia: A suite for analyzing large language models across training and scaling.
\textit{Proceedings of ICML 2023}.

\bibitem{radford2019language}
Radford, A., et al.\ (2019).
Language models are unsupervised multitask learners.
\textit{OpenAI Blog}, 1(8).

\end{thebibliography}

\end{document}
